\chapter{中央苏区的炼成}

\section{中央苏区的成长}

1930年代，中央苏区成为全国各苏区中的“中央”，绝非浪得虚名。由于朱毛红军在赣南一带的出色表现，早在1930年前后，这里已经成为中共中央和共产国际关于中国革命道路问题讨论的重心。1930年7月底，共产国际执行委员会致信中共中央，将建立巩固的农村根据地问题作为开展农村武装斗争的必要条件，强调：“完全掌握农民武装斗争的一切形式，尤其是直接着手组建惟一能保证我们巩固胜利的正规部队，只有在牢牢占领并保持具有巩固和进一步扩大苏维埃政权的足够政治经济前提的根据地的条件下才有可能”；同时指出：“我们越是迅速地具有这样的根据地，越是迅速到把武装斗争从各种独特的游击战变为正规军作战形式，我们就越能迅速地保证从组织上掌握农民革命运动，就越能迅速地保证无产阶级对农民的领导，从而保证革命的胜利。”\footnote{《共产国际、联共（布）与中国革命档案资料丛书》第7册，中央文献出版社，2002，第242页。}根据这一判断，8月8日，共产国际再电中共中央，进一步指出：

\begin{quote}
必须选择和开辟能保证组建和加强这种军队的根据地。对根据地的基本要求是：相当程度是农民运动，从容组建的可能性，获得武器的前景和保证今后能夺取一个有足够工人居民的大的行政政治中心的发展前景。目前显然赣南、闽南、粤东北地区首先能够成为这样的根据地。\footnote{《共产国际执行委员会致中共中央电（1930年8月8日）》，《共产国际、联共（布）与中国革命档案资料丛书》第9册，第189页。}
\end{quote}

这里提到的三个地区此时只有赣南地区已经有大规模的红军活动，而先后在这一地区活动过的朱德、毛泽东、彭德怀率领的部队，正是中共武装的精英。当时共产国际代表判断，在所有红军部队中，“朱德和毛泽东的军以及处在他们影响下的两个军（第3军和第12军）是最好的。彭德怀的军以及在他影响下的两个军（第8军和第16军），与他们差别不大”。\footnote{《共产国际、联共（布）与中国革命档案资料丛书》第9册，第208页。}中共中央也明确肯定：“四军是中国红军主力的主力”。\footnote{《中央给四军前委的指示信（1930年4月3日）》，《中共中央文件选集》第6册，第63页。}显然，已有根据地基础，又拥有优良红军和坚强领导人的赣南更有可能满足共产国际建立中心根据地的要求。不过考虑到自己远在千里之外，毕竟对当地情况不太熟悉，共产国际并不想在这样的具体问题上遽下定论，在有关电报后面不忘补充说：“更详细地核实这一情况只能在当地进行”。\footnote{《共产国际、联共（布）与中国革命档案资料丛书》第9册，第189页。}这是1927年莫斯科在中国遭遇挫败后，指导方式开始调整的一个例证。

共产国际提出这一计划，正是此时实际控制中共中央的李立三错误向着最为极端的方向发展之时，某种程度上，这本身就是为着减消李立三的错误而提出的。因此，在被狂热中的中共中央冷藏了近一个月后，随着进攻长沙的失败及李立三总暴动方针的破产，尤其是1930年9月中共六届三中全会召开，共产国际这一指示开始得到更多重视。首先是关于建立苏区中央局的意见得到落实。6、7月间，共产国际数次要求中共中央在苏区“成立有权威的中央局，采取一切措施尽可能加强红军”，\footnote{《共产国际执行委员会致中共中央电（1930年6月19日）》，《共产国际、联共（布）与中国革命档案资料丛书》第9册，第199页。}但正沉浸在革命高潮美好想象中的中共中央当时对此无暇顾及。周恩来同年底谈道，这时的中共中央“对于苏维埃区域的工作，并不放在党的工作之第一位……同时反对我在国际提出的建立根据地的意见，认为是右倾，保守观念。六月间，国际来电提出建立苏区的中央政府，也未被重视”。\footnote{周恩来：《关于传达国际决议的报告（1930年12月1日）》，《中共中央文件选集》第6册，第525页。}8月底，李立三的狂热渐现退潮之势后，根据共产国际指示，周恩来在中共中央总行委主席团会议上提出成立苏区中央局的建议。10月17日，中共中央政治局会议最终确定了中共苏区中央局成员名单。同时，周恩来“在政治局提出自己前往苏区的建议”，由于大家都认为“他在政治局里简直是改进党的工作和改造党不可替代的人物”，\footnote{《共产国际、联共（布）与中国革命档案资料丛书》第9册，第324页。}这一提议未能通过，但体现了此时中共中央开始高度重视苏区建设。接着，中共中央初步明确了中央根据地的范围，10月24日，中共中央政治局在关于苏区的工作计划中谈道：“我们现在确定湘鄂赣联接到赣西南为一大区域，要巩固和发展它成为苏区的中央根据地。环绕着它的首先是赣东北与湘鄂边两个苏区根据地，再则，鄂东北与闽粤赣两个苏区也很重要。”\footnote{《中央政治局关于苏维埃区域目前工作计划（1930年10月24日）》，《中共中央文件选集》第6册，第428页。}决定：“在中央苏区立即设立中央局，目的在指导整个苏维埃区域之党的组织，同时，并在苏区成立中央军事委员会以统一各苏区的军事指挥。”\footnote{《中央政治局关于苏维埃区域目前工作计划（1930年10月24日）》，《中共中央文件选集》第6册，第428页。}这是我们看到的中共中央文件中首次提出中央苏区这一概念。此后，随着形势的发展，文件中提到的赣南和闽粤赣（即闽西）两块根据地实际成为中央苏区的组成部分。1931年2月，共产国际代表使用了“江西的主要根据地”\footnote{《共产国际、联共（布）与中国革命档案资料丛书》第10册，第63页。}的提法，3月，共产国际远东局的报告中出现“朱［德］—毛［泽东］的中央（苏）区”\footnote{《共产国际、联共（布）与中国革命档案资料丛书》第10册，第184页。}这一概念，中央苏区的地位已经初步奠定。

赣南及赣西南苏区受到共产国际和中共中央的高度重视，当然和这里卓有成效的苏区建设密不可分，毛泽东、朱德、彭德怀等军政领导人在其中起着不可替代的作用。而随着苏维埃区域的持续壮大，共产国际和中共中央的关注不断加强，这里作为全国苏维埃中心的地位也更加巩固。1931年11月，苏维埃临时中央政府在江西瑞金成立，由此，赣南、闽西作为中央苏区的地位终于实至名归。

随着共产国际和中共中央日益重视根据地建设，一大批军政领导人陆续进入中央苏区，中央苏区更是人员流向的重点。1930年9月到1931年4月，中共中央政治局共向中央苏区派出67人，其中57人到达。\footnote{《共产国际、联共（布）与中国革命档案资料丛书》第10册，第289页。}苏区中央局也正式在中央苏区开始运行。1931年底，周恩来进入中央苏区。1932年下半年，共产国际东方地区书记处鉴于中共中央机关在上海难以立足，建议将中共中央迁往中央苏区。\footnote{《共产国际、联共（布）与中国革命档案资料丛书》第13册，第230页。}11月，王明致信联共（布）驻共产国际代表团，正式提出“将中央迁往中央苏区”。\footnote{《共产国际、联共（布）与中国革命档案资料丛书》第13册，第246页。}年底，共产国际决定：“采纳王明同志的建议，将中共中央、共青团中央和赤色工会总理事会从上海迁往苏区。”\footnote{《共产国际、联共（布）与中国革命档案资料丛书》第13册，第274页。}1933年初，中共中央迁往中央苏区，中央苏区成为中共中央的指挥中心。

大批领导人进入中央苏区，一时间使这里人才济济，并使苏区原有的组织体系发生巨大变化。虽然共产国际对毛泽东、朱德等苏区创建者予以高度评价，指示“对于毛泽东，必须采取最大限度的克制态度和施行同志式的影响，为他提供充分的机会在中央或中央局领导下担任负责工作”。\footnote{《共产国际、联共（布）与中国革命档案资料丛书》第13册，第352页。}中共中央则强调：“我们坚主采取一切方法，根据党的路线，缩小争论；无严重破坏纪律之事，则绝不应采取任何组织结论。”\footnote{《中共中央关于争取革命在一省与数省首先胜利的决议（1932年1月9日）》，《中共中央文件选集》第8册，第44页。}而像周恩来这样的新领导人对毛泽东也相当尊重，1932年6月，他向中共中央报告：“毛泽东身体极弱，他仍留在高山地区工作，他失眠，胃口也不好。但他和部队一起活动，在主持作战行动时精力充沛，富有才华。”\footnote{《周恩来致中央局电（1932年6月13日）》，《周恩来军事文选》第1卷，第134页。}但是一批党内现有地位高于毛泽东、朱德等领导人的到来，使毛泽东不再可能像从前那样成为苏区事实上的掌控者，而越来越被有意无意地置于边缘位置。尤其当中共中央把毛泽东视作右倾方针的代表时，不仅是他本人，他的工作作风、思想方式、应对办法都遭到质疑，与其有关的一批富有实际经验的苏区原有干部也被冷落甚而被斗争，这对苏区长远建设、发展并不有利。

历史运行真是十分复杂，正因为毛泽东、朱德等的出色表现，使赣南、闽西苏区成为中共活动的中心，而这种中心的到来，却又不可避免地影响到毛泽东等在苏区的权威。对于毛泽东而言，确实难免会有情何以堪的感觉。当然，也不能简单地认为中共中央就是摘桃派，中共中央实际上的接管中央苏区，首先是由中共严格的组织纪律所决定，在这方面，个人或团体没有讨价还价的余地；其次，各个苏区的发展固然包含着领导人的天才创造，但毕竟是在中共中央统一领导之下，与中共中央的指导和帮助不能分开；另外，中共中央到来前，中央苏区在肃AB团等问题上犯下严重错误，领导威信其实已经受到影响，黄克诚回忆其与红军干部何笃才间的一段经历，颇值参考：

\begin{quote}
何笃才是有功的干部之一。但由于他在古田会议之前朱、毛的争论当中，反对毛泽东的正确意见，从此便不受重用……他曾对我说过，毛泽东这个人很了不起！论本事，还没有一个人能超过毛泽东，论政治主张，毛泽东的政治主张毫无疑问是最正确的。我问他：既然如此，你为什么要站到反对毛泽东的一边呢？他说，他不反对毛泽东的政治路线，而是反对毛泽东的组织路线。我说：政治路线正确，组织路线有点偏差关系不大吧？他说：不行！政治路线、组织路线都不应该有偏差，都是左不得，右不得的。我问他：毛泽东的组织路线究竟有什么问题？他说：毛泽东过于信用顺从自己的人。对待不同意见的人不能一视同仁，不及朱老总宽厚坦诚……对何笃才的这番话，我是在一年以后才品味出其中的某些道理。本来，毛泽东同志在中央革命根据地军民中，已经有了很高的威望，大家都公认他的政治、军事路线正确。然而，临时中央从上海进入中央苏区后，轻而易举地夺了毛泽东的权，以错误的政治、军事路线，代替了正确的政治、军事路线。其所以会如此，苏区的同志相信党中央固然是一个重要原因。但是，如果不是毛泽东在组织路线上失掉了一部分人心，要想在中央苏区排斥毛泽东，当不会是一件容易的事情。\footnote{黄克诚：《关于毛主席评价等问题对中央提的一些意见》，《黄克诚自述》，人民出版社，1994，第142～143页。}
\end{quote}

除了新中央与老干部之间的冲突外，中共中央到来后，即就活跃于第一线的领导人而言，当时强有力的领导核心也远未形成，领导者之间、前方与后方之间的掣肘现象常有发生。周恩来在指挥第四次反“围剿”时就抱怨：“关于行动部署，尤其是许多关联到战术上问题的部署，请求中央、中央局须给前方活动以机断余地和应有的职权，否则命令我们攻击某城而非以训令指示方针，则我们处在情况变化或不利的条件下，使负责者非常困难处置。”\footnote{《周恩来、王稼祥、朱德致中央局电（1933年1月27日）》，《周恩来军事文选》第1卷，第216页。}即使是被认为同处留苏阵营的张闻天和博古，暗中也不是没有相互角力，张闻天后来曾谈到他和博古此一时期的分歧，认为：“他曾经写了一篇文章《关于苏维埃的经济政策》，暗中是驳我的个别意见的。”\footnote{张闻天：《从福建事变到遵义会议》，《遵义会议文献》，人民出版社，1985，第78页。}广昌战役后，围绕着战役方针问题，中共中央内部也有争论。张闻天回忆，广昌战役后的一次会议上，他曾提出质疑：“我批评广昌战斗同敌人死拼，遭受不应有的损失，是不对的。他（指博古——引者注）批评我，说这是普列哈诺夫反对1905年俄国工人武装暴动的机会主义思想。我当时批驳了他的这种污蔑，坚持了我的意见，结果大家不欢而散。其他到会同志，没有一个表示意见。”\footnote{张闻天：《从福建事变到遵义会议》，《遵义会议文献》，第78页。}杨尚昆后来也谈到了这场争论，他回忆的会议结果是“恩来同志当场调停，宣布散会”。\footnote{杨尚昆：《回忆长征》，《杨尚昆回忆录》，中央文献出版社，2001，第106页。}

当然，在第五次反“围剿”开始前，这些问题还不明显，而且苏区正处于其发展的高峰期。1933年的第四次反“围剿”在周恩来、朱德指挥下，首次在红军中运用大兵团伏击战法，取得重大胜利，毛泽东也不吝美言，赞誉其为“空前光荣伟大胜利”。\footnote{毛泽东：《第四次反“围剿”的伟大胜利》，《毛泽东军事文集》第1卷，第278页。}第四次反“围剿”胜利后，国民党军在江西处于守势，红军乘机在赣中、赣东北一带积极活动，扩大苏区范围，中央苏区疆域进一步扩大，跨有闽赣两省数十个县。到第五次反“围剿”前夕，中央苏区在江西达到极盛，除占据赣南一半以上地域外，北延到南城、黎川地区，面积在4万平方公里左右。

随着中央苏区的不断壮大，人才空前集聚并拥有了中共中央机关的中央苏区领导层当然不是毫无作为。七大的建军报告初稿中曾写道：“事实上以毛泽东同志为首的党的正确路线，本身是一个发展的东西一个发展的过程，是在不断克服困难和错误之中把自己坚强起来，我们更应该坚持这种服从真理追求真理的正确态度。”\footnote{《关于建军问题的报告初稿（1945年）》，转见杨奎松《“中间地带”的革命》，第340页。}这一判断在苏区时期其实同样适应。也就是说，虽然毛泽东在苏区发展壮大中表现的卓越才能超乎于包括博古在内的中共中央许多人之上，但不能认为苏区建设就一定会由于毛泽东地位的变化而升降。中共中央进入苏区后取得了一些值得列举的政绩：1933年，在连续两年歉收后，中央苏区迎来了苏维埃政府建立后的第一个丰收年，产量接近革命前的水平，毛泽东1933年对才溪乡的调查展示的虽是比较好的状况，但仍有一定的代表性：“暴动后（一九二九至一九三一年），生产低落约百分之二十。一九三二年恢复了百分之十，今年（一九三三）比去年增加二成（杂粮如番薯、豆子、芋子、大薯等，则比去年增加了百分之五十），超过了暴动前百分之十。暴动后全区荒了许多田，去年开发了一小部分。今年大开，开了一千三百多担。”\footnote{毛泽东：《才溪乡调查》，《毛泽东农村调查文集》，人民出版社，1982，第344页。}虽然收成的好坏和自然气候有着一定的关系，但土地明确归农户所有等一系列刺激农民生产积极性的措施在其中起到的作用仍不可忽视。同时，中共中央在经济、劳动等政策上也作出了一系列比较符合实际的决策。军队和地方的正规化建设由于中共中央的到来都有了长足的进展。第四次反“围剿”时，由于对军队基本战斗技能培训明显加强，红军将领深切感受到：“射击技能之不够在过去成为部队普遍的现象，而这次战斗中射击技能之进步并成绩确有不可抹煞之事实。”\footnote{《红一方面军第四次反“围剿”战役总结（1933年4月）》，《朱德军事文选》，第134页。}

中共中央到来后，随着大批具有相当文化水准人员的陆续进入，中央苏区的文化教育空前发达。各种杂志、报纸纷纷创刊，反映着当时中共中央的态度，舆论批评的气氛尚较活跃，并由此给后人留下了更完整、全面认识中央苏区的珍贵史料。在肃反问题上，虽然肃反扩大化趋势仍不能克服，但富田事变前后的严重肃反错误已被纠正，党内的所谓“残酷斗争”主要体现在思想方面，组织上的措施明显比此前谨慎，即使是被集中批判的对象罗明，也没有遭受肉体上的摧残。大批干部的到来也使监督体系和民主决策有了一定的发展。军队正规化建设正在加强，干部的使用更程序化，苏维埃政权的组织体制有了相当程度的提高。这些，代表着苏维埃政权军事、政治上不断发展的现实和趋势，中共中央作为工作的主持者在其中发挥的作用不应完全抹杀。所以，1944年六届七中全会决议中明确指出：“自四中全会至遵义会议期间，党中央的领导路线是错误的，但尚有其正确的部分，应该进行适当的分析，不要否定一切。”\footnote{毛泽东：《关于若干历史问题的决议》，《毛泽东选集》第3卷，第955页。}这一说法是中肯的。

当然，和毛泽东时代注重实际、认真调查研究、在实践中艰辛探索从而迅速使苏区获得发展壮大相比，毛泽东逐渐淡出决策层后，中央苏区确实存在着一些严重的问题，毛泽东曾谈道：“在1927～1935年而特别是1931～1935年时期我党曾经因为政策过左而陷于孤立，处于极端危险的地位。”\footnote{毛泽东：《关于陕甘宁边区的文化教育问题》，《毛泽东文集》第3卷，第122页。}他所作的更具体解释是：

\begin{quote}
土地革命战争时期，党内机会主义的主要特点是“左”……把赤白对立绝对化；对中小资产阶级实行过左政策；片面强调工人利益而把工商业很快搞垮了；主张地主不分田、富农分坏田，并损伤了一部分中农的利益。当然，我们党在农村中还是有群众的，不能说是在农民中完全孤立。总之，土地革命战争时期实行“左”的政策的结果，我们没能孤立蒋介石，而是孤立了自己。\footnote{毛泽东：《对〈关于若干历史问题的决议〉草案的说明》，《毛泽东文集》第3卷，第287页。}
\end{quote}

\section{新形势与新任务}

这样一个发展与问题共存、期待与挫折并进中的中央苏区，在1933年秋季面临着新的形势和变化。1933年1月30日，刚刚抵达苏区的张闻天发表文章阐述当前政局，指出：“自从苏维埃与红军在粉碎帝国主义国民党的三次‘围剿’后在中国广大领土上威胁着帝国主义国民党的统治，苏维埃临时中央政府宣告成立以来……国民党的统治更形动摇，中国革命危机进一步的深入，全国反帝反国民党的斗争更带着群众性，无产阶级与农民及一切劳动民众的团结更加巩固，更加表示了他们对于中国共产党与中华苏维埃共和国临时中央政府的拥护。”\footnote{洛甫：《中国事变与中国共产党》，《斗争》第1期，1933年2月4日。}

1933年，进入中央苏区的中共中央确实看到他们面前呈现的是一幅相当积极的画面。第四次反“围剿”胜利后，中央苏区形势向着更好的方向发展，苏区版图进一步扩大。当时中央苏区周边的形势是：浙赣边的新红十军总人数不下6000人，在不断扩充之中；新湘鄂赣的红十六军总人数2000余人，仍能保存其实力；湘赣边的红八军总人数约三个师，能在永新、莲花、宁冈、茶陵一带与陈光中部63师等对抗，并不断保持其游击区；湘鄂苏区的红三军，亦在鹤峰、桑植一带巩固、扩大根据地。\footnote{《共产国际、联共（布）与中国革命档案资料丛书》第13册，第386页。}在全国各苏区中，中央苏区的优势十分明显。受此鼓舞，中央苏区领导人进一步提出了扩红百万的庞大计划，并准备将中央苏区向福建方向延伸。\footnote{《共产国际、联共（布）与中国革命档案资料丛书》第13册，第425～427页。}

1933年5月，蒋介石坐镇南昌，开始布置对苏区的第五次“围剿”。但在战争之初，国民党军并未立即取得明显的优势，中共依然对战争前景抱持乐观态度。在共产国际派来的军事顾问李德（Otto Braun）到达中央苏区之前，前方周恩来、朱德和后方的博古等对以阵地战、堡垒战对堡垒战的观点基本是一致的。\footnote{参见《共产国际、联共（布）与中国革命档案资料丛书》第14册，第63、67页。}这种一致体现着中共中央到来后党内认识的趋同，也反映着当时大多数中共领导人的认知水平。但是，面对国民党军大兵压境，这种以硬碰硬的战略的选择，对处于弱势的中共其实是相当不利的。蒋介石的德国顾问冯·塞克特（von Seeckt）对他的信心来源作了清晰的解释：“我们不应该将我们的部队投入到旷野之中，我们应该用碉堡政策来窒息敌人的生命力。”\footnote{Otto Braun, \textit{A Comintern Agent in China}, 1932--1939 (Stanford: Stanford University Press, 1982), p. 57.}这就是以“乌龟壳”战术对中共之“短促突击”：步步为营、稳扎稳打，将红军引向国民党军所希望的阵地战、消耗战。

尽管如此，面对中共在赣南闽西的发展壮大，当时多数人并没有立即陷入悲观。因为红军虽然仍保持着游击战的传统，但也具有了正规战的能力和实力。周恩来曾乐观地估计：“由于我们没有利用时机在军事上大踏步地前进……我们没有能够在有利时机将国民党军消灭在个别方向上。尽管如此，我们仍可能夺回主动权。”\footnote{Otto Braun, \textit{A Comintern Agent in China}, p. 56.}

\section{从毛泽东眼中走进苏区}

在展开论述国共双方的这场殊死角逐之前，我们有必要进一步深入到中央苏区的内部，对当年毛泽东所说的苏维埃革命根据地的社会状况作出描绘。恰在第五次“围剿”前夕，毛泽东对江西兴国长冈乡、福建上杭才溪乡作了详细调查。这两份调查报告在中共党内受到高度重视，被作为典型材料印发全军，对研究中央苏区的社会实况极具价值。毛泽东的调查是在他的政治生涯逐渐滑向低谷之时进行的。在此前后，毛泽东还做了寻乌调查、木口村调查等，这些调查为理解当年苏维埃区域的社会状况提供了宝贵而翔实的材料。

从地理形势看，中央苏区所在的赣南、闽西地区地处赣江、汀江上游，武夷山纵贯其间，海拔多在1000米左右。区域内山岭重叠，交通不便，经济落后。然而，正是这种地理条件为红军开展游击战争提供了天然屏障。毛泽东指出，“有了农村之后，我们必须注意经济工作……这是一个伟大的任务，一个伟大的阶级斗争”。\footnote{毛泽东：《必须注意经济工作》，《毛泽东选集》第1卷，第119页。}事实上，毛泽东所做的调查详尽考察了苏区农村的社会结构和阶级状况，为我们留下了一份珍贵的社会记录。

毛泽东等人的调查显示，赣南、闽西农村社会结构和全国大部分地区一样，宗族势力在当地具有举足轻重的影响力。赣南、闽西地区宗族聚居现象十分普遍，宗族不仅拥有社会政治权威，而且还占有土地——公田。公田和宗族势力的存在，一方面使当地社会呈现出某种以血缘为纽带的共同体特征，另一方面也造成社会结构的相对封闭和保守。

基于对苏区社会结构、阶级关系的深入了解，毛泽东后来在《中国革命战争的战略问题》等著作中，反复强调红军力量的有限性，以及必须从这个实际出发决定红军的战略战术。毛泽东的这些思想，和他在中央苏区所做的深入调查研究是分不开的。

\subsection{宗族势力的影响}

赣南、闽西地区宗族聚居现象十分普遍，单姓村和主姓村占有相当比例。根据毛泽东在寻乌的调查，该县共有约2.68万户居民，分属200余姓，其中刘、陈、潘、谢、何等大姓占有绝对优势。\footnote{毛泽东：《寻乌调查》，《毛泽东农村调查文集》，第68页。}“聚族而居，族必有祠”是这一地区的重要特征。祠堂不仅是宗族祭祀的场所，也是宗族议事、执法的中心。宗族通过族规、族训对族人实施有效的控制和管理，形成了一套相对完整的社会自治系统。

宗族力量的强大还体现在经济方面，主要形式是拥有公田。毛泽东在寻乌调查中发现：“寻乌公田多，成了普遍现象。”公田名义上为宗族共有，实际上为宗族中的士绅所控制。不过由于公田的存在，一些贫困的宗族成员可以通过租种公田维持生计，这在客观上形成了一定的社会缓冲机制，也使得阶级矛盾在一定程度上被宗族关系所掩盖。中共在进入这一地区后，很快认识到了宗族问题的复杂性，提出了打击宗族势力、没收公田的口号，但由于宗族观念深入人心，实际推行中遇到了不少阻力。

\subsection{中农问题}

在有关1920至1940年代中国农村社会的描述中，地主、富农和贫雇农作为社会对立的两极，往往受到充分的重视，而在这两极间的一个庞大社会阶层——中农却常常遭到忽视。其实，当时的调查材料反映出中农在农村中具有不可替代的地位和影响。

由于中农是作为一种阶级分析的方法而提出的名词，1949年前，除受马克思主义影响得出的调查成果外，大部分调查并不使用中农这一概念，而是以地主、自耕农、佃农进行种类划分，即使受中共影响进行的调查，对中农的定义也不尽相同，因此中农在人口中的比例和社会中的影响难以得到准确反映。土地革命广泛开展后，随着中共积累起更多阶级划分经验，中农形象渐渐鲜明，毛泽东给中农下的定义是：

\begin{quote}
中农许多都占有土地。有些中农只占有一部分土地，另租入一部分土地。有些中农并无土地，全部土地都是租入的。中农自己都有相当的工具。中农的生活来源全靠自己劳动，或主要靠自己劳动。中农一般不剥削别人，许多中农还要受别人小部分地租债利等剥削。但中农一般不出卖劳动力。另一部分中农（富裕中农）则对别人有轻微的剥削，但非经常的和主要的。\footnote{毛泽东：《怎样分析农村阶级》，《毛泽东选集》第1卷，第128页。}
\end{quote}

对于强调阶级分析的中共党人而言，中农的划分具有相当重要的意义，也是研究者不能忽视的话语。

从采用中农概念的多项调查提供的数据看，中农人口和占地比例在农村通常是最大的阶层，一般达到当地农村的30\%左右。从人口比例看，中农在农村各阶层中比例一般仅次于贫农。中农在农村人口中占有的重要地位，使其成为影响农村的一个不可忽视的因素。

作为社会经济发展水平的标志之一，一般而言，在经济相对发达地区，中农比例较高，而在相对贫穷、落后地区，中农比例要低一些。苏区由于多处山区，经济相对落后，中农发育受到影响，较之社会平均指标要低一些，但即使如此，中农仍占相当比例。即便在苏区，中农人口仍占到1/4左右。

从占地比例看，中农占地一般稍稍超过人口比例达到35\%左右。从生产资料主要是耕牛占有情况看，中农更具有较大优势。据对福建14县22个村的调查，耕牛占有情况是：中农占有耕牛总数的55.69\%，贫雇农占有32.63\%，富农占有8.76\%，地主占有1.12\%，其他阶层占有1.8\%。\footnote{《福建省农村调查》，华东军政委员会土地改革委员会编印，1952，第189页。}中农在人口、土地及生产资料占有上的重要地位，使其动向对农村社会具有不可忽视的影响。

苏维埃革命时期的中共话语中，中农一直是联合的对象，只是在联合的地位和紧密度上随着判断的变化有所差异。随着苏维埃革命的深入进行，中农在农民革命中的地位得到更充分的认识，中共六大决议已经把联合中农放到战略高度予以考量，强调：

\begin{quote}
联合中农是保证土地革命胜利的主要条件。贫农与农村无产阶级在工人阶级领导之下而斗争，是土地革命的主要动力，而与中农联合是保证土地革命胜利的主要条件，中国共产党提出之没收一切地主土地分配给无地或少地的农民的政纲，必能得到广大的中农群众的拥护。因为中农群众也是受地主阶级各种封建剥削压迫群众中的一部份。\footnote{《中共第六次全国代表大会关于农民问题决议案》，《第一第二次国内革命战争时期土地斗争史料选编》，人民出版社，1981，第245页。}
\end{quote}

虽然就理论上而言，中共从未怀疑过中农的重要性，但中农和贫雇农毕竟是具有不同经济利益和经济地位的两个群体，尤其在动荡的政治运动期间，这种利益冲突更容易被诱发。作为一个马克思主义政党，从阶级分析立场出发，中共对贫雇农的信任度当然要高于中农。在此背景下，像贫农团这样被认为是“乡村苏维埃政权的柱石”的贫雇农的群众组织，在阶级话语中拥有绝对强势地位。

当然，中共对中农的重视确实是真诚的，尽力使中农成为苏区社会的主流阶层，符合中共的战略利益。张闻天1933年更从动态角度具体阐发了制定正确的中农政策的意义：“中农在中国苏维埃区域内占最大数量。在许多区域内虽是贫农出身的农民占多数，但是在从土地革命之后，这些贫农将逐渐的变为中农。这一特点，对于我们有极重大的意义。如若我们不能巩固同中农的联合，苏维埃革命是没法胜利的。”张闻天明确指出：“在开展农村中的阶级斗争，我们不要一刻忘记联合中农，这需要我们党的领导同志，最审慎的细心来感觉它们每一个的情绪与他们的要求。”\footnote{张闻天：《苏维埃政权下的阶级斗争》，《斗争》第15期，1933年6月15日。}

\subsection{公田问题}

中国传统社会宗族势力发达，在许多地区，宗族不仅拥有社会政治权威，而且占有土地。中央苏区所在的赣南、闽西地区，主要以宗族占有形式出现的公田占据很大比例。根据1950年初华东军政委员会的统计，浙江公堂土地占比为16.35\%，安徽为4.17\%，苏南为5.9\%。\footnote{《华东区土地改革成果统计》，华东军政委员会土地改革委员会编印，1952，第4～5页。}中央苏区中心区的江西属公田发达地区，国共两党的调查中都不约而同地注意到这里占地广大的公田，中共江西省委1932年5月的工作报告记载：“江西公堂祠堂的土地特别多。”\footnote{《江西苏区中共省委工作总结报告（1932年5月）》，《中央革命根据地史料选编》（上），第445页。}国民党方面在武力恢复对苏区控制后，也报告这里“公田甚多”。\footnote{《陈诚等呈收复地土问题亟待解决请迅予布告》，《军政旬刊》第32期，1934年8月31日。}

公田的大量存在，对中共领导的土地革命产生了双重影响。一方面，公田的分配触及的利益面相对较小，不会立即引起私人地主的大规模反抗，为中共开展土地革命提供了方便的资源；另一方面，由于公田与宗族制度的紧密联系，没收和分配公田不可避免地触动了传统的宗族秩序，引发了社会关系的剧烈变动。

\section{国共较量中的地缘政治}

从更广阔的视野看，中央苏区所在的赣南、闽西之所以在1930年代成为全国苏维埃运动的中心，除了中共在此地的不懈努力外，和这里独特的地缘政治环境也有着密切的关系。当年毛泽东在《中国的红色政权为什么能够存在？》和《井冈山的斗争》等著作中，从理论上阐述了这一地区成为革命中心的原因。他认为，红色政权存在和发展需要两个基本条件：一是白色政权的长期分裂与战争，二是受过民主革命影响的地区。\footnote{毛泽东：《中国的红色政权为什么能够存在？》，《毛泽东选集》第1卷，第48～49页。}

赣南、闽西地区恰好完全符合这两个条件。从第一个条件看，江西、福建两省是国民政府中央统治力量及地方政治军事势力都相当薄弱的地区。南京国民政府成立后，江西控制权长期处于客籍军人手中，与中央政府若即若离，和地方也是各怀心思；福建则是民军蜂起，各不相下，省政几成瘫痪状态。由于中央权威软弱，地方力量又极不发展，当中共在赣东北展开革命宣传时，地方政权十分惊恐，甚至不得不采取放任态度：“以前他们所张的反共标语，县长下令取消了，他说：‘共党是惹不得的，越惹越厉害，到是不管的好’。”\footnote{《鄱阳党团工作报告（1927年11月）》，《闽浙皖赣革命根据地》，中共党史出版社，1991，第58页。}因此，中共在这里的发展具有得天独厚的条件。

从第二个条件看，国民革命曾经在江西、福建掀起的巨澜，为中共在两省的组织发展提供了十分有利的基础。因此，当国共合作破裂，中共独立开展苏维埃革命时，其中心地区主要围绕着国民革命基本区展开绝非偶然。统治力量的薄弱、大山屏蔽的自然环境、国民革命运动打下的良好基础、赣南闽西背靠广东这一与南京政府保持半独立状态地区的特殊地理态势，为红军和苏维埃的发展提供了难得的有利条件，也是这里成为苏维埃革命中心区的主要原因。

强调赣南闽西成为苏维埃革命中心的环境、力量因素，绝不意味着否认这里存在内在的革命动力。事实上，在农村贫困的背景下，这一要求在全国普遍存在，关键是，其是否能被调动和发挥。所以，虽然我们在江西、福建看到并不是十分畸形的地权关系，但并不影响这里成为革命的中心；而地主、富农在农村经济危机下遭遇的困境，也不能使他们免于革命的打击。

当我们顺着毛泽东的调查走进中央苏区后，为尊重历史的客观性，不妨再听听对手方的看法。陈诚作为当年“剿共”的悍将、日后国民党方面的重臣，他的看法在国民党方面应该分量不轻，其对赣南成为中共根据地的分析更不无见地，某种程度上恰可与毛泽东当年的分析互为印证，他的分析很长，但颇值一读，引录如下：

\begin{quote}
第一因为地理环境关系，赣南位于赣江上游，地势高峻，山岭重迭，交通极为不便，这是打出没无定的游击战最理想的地带。共党最擅长的就是打游击战，所以他们选定了赣南作主要根据地。而且赣南的经济条件也很优越。赣南虽然山多，但因侵蚀年久，山间溪谷，多冲积成局部平原，颇适宜于耕种。前章提到的瑞金，就是“种一年吃三年”的好地方。其他各县虽不都和瑞金一样，可是出产的种类数量，都很丰富，维持一个经济生活自给自足的局面，是可能的，所以他们就看中了赣南。
\end{quote}

\begin{quote}
第二因为政治环境关系。江西政治环境最利于共党发展，其故有二：一、江西东面的福建，十九路军驻入以前，政府于此素乏经营，十九路军驻入以后，即逐渐反动，为政府之患。江西南面的广东，形同割据，反抗中央，固已匪伊朝夕。江西西面的湖南，与政府同床异梦，于共党亦无所害。故共党据赣南，所虑者惟北面耳。二、民国以来，江西遭受军阀的摧残，为各省之冠。北伐成功后，人民对于改善政治环境的要求很高，希望非常之大。不想当时国家统一徒俱虚名，军阀割据，内乱迭起，政府对于改善地方政治，有心无力，赣南山乡辽远，遂致更成化外。人民的希望破灭了，在艰苦中挣扎生活，似乎毫无出头之日。这种环境，是共党最易欺骗民众的。所谓贫穷是繁殖红色细菌的温床，就是这个意思。这是共党看中了赣南的另一原因。
\end{quote}
